% Options for packages loaded elsewhere
\PassOptionsToPackage{unicode}{hyperref}
\PassOptionsToPackage{hyphens}{url}
%
\documentclass[
  ignorenonframetext,
  aspectratio=43,
]{beamer}
\usepackage{pgfpages}
\setbeamertemplate{caption}[numbered]
\setbeamertemplate{caption label separator}{: }
\setbeamercolor{caption name}{fg=normal text.fg}
\beamertemplatenavigationsymbolshorizontal
% Prevent slide breaks in the middle of a paragraph
\widowpenalties 1 10000
\raggedbottom
\setbeamertemplate{part page}{
  \centering
  \begin{beamercolorbox}[sep=16pt,center]{part title}
    \usebeamerfont{part title}\insertpart\par
  \end{beamercolorbox}
}
\setbeamertemplate{section page}{
  \centering
  \begin{beamercolorbox}[sep=12pt,center]{part title}
    \usebeamerfont{section title}\insertsection\par
  \end{beamercolorbox}
}
\setbeamertemplate{subsection page}{
  \centering
  \begin{beamercolorbox}[sep=8pt,center]{part title}
    \usebeamerfont{subsection title}\insertsubsection\par
  \end{beamercolorbox}
}
\AtBeginPart{
  \frame{\partpage}
}
\AtBeginSection{
  \ifbibliography
  \else
    \frame{\sectionpage}
  \fi
}
\AtBeginSubsection{
  \frame{\subsectionpage}
}

\usepackage{amsmath,amssymb}
\usepackage{lmodern}
\usepackage{iftex}
\ifPDFTeX
  \usepackage[T1]{fontenc}
  \usepackage[utf8]{inputenc}
  \usepackage{textcomp} % provide euro and other symbols
\else % if luatex or xetex
  \usepackage{unicode-math}
  \defaultfontfeatures{Scale=MatchLowercase}
  \defaultfontfeatures[\rmfamily]{Ligatures=TeX,Scale=1}
\fi
\usetheme[]{default}
\usecolortheme{default}
% Use upquote if available, for straight quotes in verbatim environments
\IfFileExists{upquote.sty}{\usepackage{upquote}}{}
\IfFileExists{microtype.sty}{% use microtype if available
  \usepackage[]{microtype}
  \UseMicrotypeSet[protrusion]{basicmath} % disable protrusion for tt fonts
}{}
\makeatletter
\@ifundefined{KOMAClassName}{% if non-KOMA class
  \IfFileExists{parskip.sty}{%
    \usepackage{parskip}
  }{% else
    \setlength{\parindent}{0pt}
    \setlength{\parskip}{6pt plus 2pt minus 1pt}}
}{% if KOMA class
  \KOMAoptions{parskip=half}}
\makeatother
\usepackage{xcolor}
\newif\ifbibliography
\setlength{\emergencystretch}{3em} % prevent overfull lines
\setcounter{secnumdepth}{-\maxdimen} % remove section numbering


\providecommand{\tightlist}{%
  \setlength{\itemsep}{0pt}\setlength{\parskip}{0pt}}\usepackage{longtable,booktabs,array}
\usepackage{calc} % for calculating minipage widths
\usepackage{caption}
% Make caption package work with longtable
\makeatletter
\def\fnum@table{\tablename~\thetable}
\makeatother
\usepackage{graphicx}
\makeatletter
\def\maxwidth{\ifdim\Gin@nat@width>\linewidth\linewidth\else\Gin@nat@width\fi}
\def\maxheight{\ifdim\Gin@nat@height>\textheight\textheight\else\Gin@nat@height\fi}
\makeatother
% Scale images if necessary, so that they will not overflow the page
% margins by default, and it is still possible to overwrite the defaults
% using explicit options in \includegraphics[width, height, ...]{}
\setkeys{Gin}{width=\maxwidth,height=\maxheight,keepaspectratio}
% Set default figure placement to htbp
\makeatletter
\def\fps@figure{htbp}
\makeatother

\usepackage{ragged2e}
\usepackage{etoolbox}

\apptocmd{\frame}{}{\justifying}{} 

\setbeamertemplate{enumerate items}[default]
\setbeamertemplate{itemize items}[circle]
\setbeamersize{text margin left=2.25em, text margin right=2.25em}

\let\OLDitemize\itemize
\renewcommand\itemize{
    \OLDitemize\addtolength{\itemsep}{3pt plus 1pt minus 1pt}
}
\let\OLDenumerate\enumerate
\renewcommand\enumerate{
    \OLDenumerate\addtolength{\itemsep}{3pt plus 1pt minus 1pt}
}

\usepackage{orcidlink}
\usepackage{amsmath}
\DeclareMathOperator*{\argmax}{argmax}

\makeatletter
\def\@listii{\leftmargin\leftmarginii
              \topsep    1ex
              \parsep    0\p@   \@plus\p@
              \itemsep   \parsep}
\makeatother

\usepackage{lipsum}

\let\olditem\item
\renewcommand\item{\olditem\justifying}

\definecolor{blue}{RGB}{0,114,178}
\definecolor{red}{RGB}{213,94,0}
\definecolor{yellow}{RGB}{240,228,66}
\definecolor{green}{RGB}{0,158,115}

\definecolor{MyBackground}{RGB}{255,253,218}

\setbeamercolor{section in toc}{fg=blue}
\setbeamercolor{subsection in toc}{fg=red}

\setbeamercolor{frametitle}{fg=blue}
\setbeamercolor{title}{fg=blue}

\setbeamertemplate{footline}[frame number]
\setbeamertemplate{navigation symbols}{} 
\setbeamertemplate{itemize items}{-}

\setbeamercolor{itemize item}{fg=blue}
\setbeamercolor{itemize subitem}{fg=blue}
\setbeamercolor{enumerate item}{fg=blue}
\setbeamercolor{enumerate subitem}{fg=blue}
\setbeamercolor{button}{bg=MyBackground,fg=blue}

\newcounter{daggerfootnote}
\newcommand*{\daggerfootnote}[1]{%
    \setcounter{daggerfootnote}{\value{footnote}}%
    \renewcommand*{\thefootnote}{\fnsymbol{footnote}}%
    \footnote[2]{#1}%
    \setcounter{footnote}{\value{daggerfootnote}}%
    \renewcommand*{\thefootnote}{\arabic{footnote}}%
}

\AtBeginSection{
    \begingroup
    \setbeamercolor{background canvas}{bg=yellow}
    \begin{frame}
        \begin{centering}
            \usebeamerfont{section title}\huge\color{black}\insertsection\par
        \end{centering}
    \end{frame}
    \endgroup
}
\makeatletter
\makeatother
\makeatletter
\makeatother
\makeatletter
\@ifpackageloaded{caption}{}{\usepackage{caption}}
\AtBeginDocument{%
\ifdefined\contentsname
  \renewcommand*\contentsname{Table of contents}
\else
  \newcommand\contentsname{Table of contents}
\fi
\ifdefined\listfigurename
  \renewcommand*\listfigurename{List of Figures}
\else
  \newcommand\listfigurename{List of Figures}
\fi
\ifdefined\listtablename
  \renewcommand*\listtablename{List of Tables}
\else
  \newcommand\listtablename{List of Tables}
\fi
\ifdefined\figurename
  \renewcommand*\figurename{Figure}
\else
  \newcommand\figurename{Figure}
\fi
\ifdefined\tablename
  \renewcommand*\tablename{Table}
\else
  \newcommand\tablename{Table}
\fi
}
\@ifpackageloaded{float}{}{\usepackage{float}}
\floatstyle{ruled}
\@ifundefined{c@chapter}{\newfloat{codelisting}{h}{lop}}{\newfloat{codelisting}{h}{lop}[chapter]}
\floatname{codelisting}{Listing}
\newcommand*\listoflistings{\listof{codelisting}{List of Listings}}
\makeatother
\makeatletter
\@ifpackageloaded{caption}{}{\usepackage{caption}}
\@ifpackageloaded{subcaption}{}{\usepackage{subcaption}}
\makeatother
\makeatletter
\@ifpackageloaded{tcolorbox}{}{\usepackage[many]{tcolorbox}}
\makeatother
\makeatletter
\@ifundefined{shadecolor}{\definecolor{shadecolor}{rgb}{.97, .97, .97}}
\makeatother
\makeatletter
\makeatother
\ifLuaTeX
  \usepackage{selnolig}  % disable illegal ligatures
\fi
\IfFileExists{bookmark.sty}{\usepackage{bookmark}}{\usepackage{hyperref}}
\IfFileExists{xurl.sty}{\usepackage{xurl}}{} % add URL line breaks if available
\urlstyle{same} % disable monospaced font for URLs
\hypersetup{
  pdftitle={Seminar 5},
  pdfauthor={Gabriel E. Cabrera Guzmán},
  hidelinks,
  pdfcreator={LaTeX via pandoc}}

\title[Seminar 5]{Seminar
5\daggerfootnote{\texttt{gabriel.cabreraguzman@postgrad.manchester.ac.uk}}}
\subtitle{Statistics 1 (Estimation) Solutions}
\author[Gabriel E. Cabrera Guzmán]{
    {Gabriel E. Cabrera
Guzmán\textsuperscript{\normalsize\orcidlink{0000-0000-0000-0000}}} \medskip
}
\date{November 21, 2025}
\institute[]{
    The University of Manchester\\
Alliance Manchester Business School \bigskip
}

\definecolor{quarto-callout-note-color}{RGB}{0,114,178}
\definecolor{quarto-callout-note-color-frame}{RGB}{0,114,178}
\definecolor{quarto-callout-warning-color}{RGB}{213,94,0}
\definecolor{quarto-callout-warning-color-frame}{RGB}{213,94,0}
\begin{document}
\frame{\titlepage}
\ifdefined\Shaded\renewenvironment{Shaded}{\begin{tcolorbox}[boxrule=0pt, breakable, enhanced, borderline west={3pt}{0pt}{shadecolor}, frame hidden, sharp corners, interior hidden]}{\end{tcolorbox}}\fi

\begin{frame}{Exercise 1}
\protect\hypertarget{exercise-1}{}
When a machine is working properly, filled chocolate bars have a weight
of 30g with a standard deviation of 1g. For quality control a sample of
50 bars is taken during a day and the average weight (\(\bar{X}\)) is
calculated. At the end of a particular day the sample average is 29.7.
Calculate a 95\% confidence interval for the mean weight of a bar on
this day.
\end{frame}

\begin{frame}{Exercise 1 (Cont'd)}
\protect\hypertarget{exercise-1-contd}{}
Since the population standard deviation is known (\(\sigma = 1\)) and
the sample size is large (\(n = 50\)), we can use the standard normal
distribution to construct the confidence interval. To find the
confidence interval, we first need to find \(Z_{\alpha/2}\). We want to
construct a 95\% confidence interval. Therefore, the confidence level is
\(1-\alpha = 0.95\), and \(\alpha = 0.05\). Considering that we have a
two-sided confidence interval, the error term in each tail must be
\(\alpha/2 = 0.025\). Thus, \(\Pr(Z < \alpha) = 0.975\). To find
suitable \(Z_{\alpha/2} = \alpha\), search for the probability of 0.975
among the numbers (probabilities) in the standard normal distribution
table. We find that \(\Pr(Z < 1.96) = 0.975\). Therefore, we have
\(Z_{\alpha/2} = 1.96\).
\end{frame}

\begin{frame}{Exercise 1 (Cont'd)}
\protect\hypertarget{exercise-1-contd-1}{}
\pause

Construct the confidence interval and replace \(\bar{x} = 29.7\),
\(Z_{\alpha/2} = 1.96\), \(n = 50\), and \(\sigma = 1\) as follows:

\pause

\[
\bar{x} - Z_{\alpha/2} \times \frac{\sigma}{\sqrt{n}} < \mu < \bar{x} + Z_{\alpha/2} \times \frac{\sigma}{\sqrt{n}}
\]

\pause

\[
29.7 - 1.96 \times \frac{1}{\sqrt{50}} < \mu < 29.7 + 1.96 \times \frac{1}{\sqrt{50}}
\]

\pause

\[
29.423 < \mu < 29.977
\]
\end{frame}

\begin{frame}{Exercise 2}
\protect\hypertarget{exercise-2}{}
A regular daily flight is due to arrival at 18:00. Its arrival time has
a normal distribution with mean 18:10 hours and standard deviation 10
minutes. A new traffic control system is installed but is criticised for
causing delays in arrivals. To investigate this, a sample of 20 flights
gave a mean arrival time of 18:12. Calculate the 95\% confidence
interval for the new mean arrival time.
\end{frame}

\begin{frame}{Exercise 2 (Cont'd)}
\protect\hypertarget{exercise-2-contd}{}
Since the population standard deviation is known (\(\sigma = 10\)) and
the population distribution is normally distributed, we can use the
standard normal distribution to construct the confidence interval. To
find the confidence interval, we first need to find \(Z_{\alpha/2}\). We
want to construct a 95\% confidence interval. Therefore, the confidence
level is \(1-\alpha = 0.95\), and \(\alpha = 0.05\). Considering that we
have a two-sided confidence interval, the error term in each tail must
be \(\alpha/2 = 0.025\). Thus, \(\Pr(Z < \alpha) = 0.975\). To find
suitable \(Z_{\alpha/2} = \alpha\), search for the probability of 0.975
among the numbers (probabilities) in the standard normal distribution
table. We find that \(\Pr(Z < 1.96) = 0.975\). Therefore, we have
\(Z_{\alpha/2} = 1.96.\)
\end{frame}

\begin{frame}{Exercise 2 (Cont'd)}
\protect\hypertarget{exercise-2-contd-1}{}
\pause

Construct the confidence interval and replace \(\bar{x} = 1812\),
\(Z_{\alpha/2} = 1.96\), \(n = 20\), and \(\sigma = 10\) as follows:

\pause

\[
\bar{x} - Z_{\alpha/2} \times \frac{\sigma}{\sqrt{n}} < \mu < \bar{x} + Z_{\alpha/2} \times \frac{\sigma}{\sqrt{n}}
\]

\pause

\[
1812 - 1.96 \times \frac{10}{\sqrt{20}} < \mu < 1812 + 1.96 \times \frac{10}{\sqrt{20}}
\]

\pause

\[
1807.617 < \mu < 1816.383
\]
\end{frame}

\begin{frame}{Exercise 3}
\protect\hypertarget{exercise-3}{}
A sample of 10 taxi fares from the station to the centre of Norwich gave
a mean of £3.80 and a sample standard deviation of 80p. Calculate the
95\% confidence interval for the mean taxi fare. What assumptions are
necessary?
\end{frame}

\begin{frame}{Exercise 3 (Cont'd)}
\protect\hypertarget{exercise-3-contd}{}
Since the population standard deviation is unknown, population
distribution is unknown, and the sample size is small. We will assume
that the population distribution is normally distributed so that we can
use Student's t distribution to construct the confidence interval. To
find the confidence interval, we first need to find \(t_{\alpha/2}\). We
want to construct a 95\% confidence interval. Therefore, the confidence
level is \(1-\alpha = 0.95\), and \(\alpha = 0.05\). Considering that we
have a two-sided confidence interval, the error term in each tail must
be \(\alpha/2 = 0.025\). Thus, \(\Pr(T < \alpha) = 0.975\). To find
suitable \(t_{\alpha/2} = \alpha\), search for the probability of 0.975
among the numbers (probabilities) with a degree of freedom of
\(n-1 = 10-1 = 9\) in the Student's t distribution table. We find that
\(\Pr(T < 2.262) = 0.975\). Therefore, we have \(t_{\alpha/2} = 2.262\).
\end{frame}

\begin{frame}{Exercise 3 (Cont'd)}
\protect\hypertarget{exercise-3-contd-1}{}
\pause

Construct the confidence interval and replace \(\bar{x} = 3.8\),
\(t_{\alpha/2} = 2.262\), \(n = 10\), and \(s = 0.8\) as follows:

\pause

\[
\bar{x} - t_{\alpha/2} \times \frac{s}{\sqrt{n}} < \mu < \bar{x} + t_{\alpha/2} \times \frac{s}{\sqrt{n}}
\]

\pause

\[
3.8 - 2.262 \times \frac{0.8}{\sqrt{10}} < \mu < 3.8 + 2.262 \times \frac{0.8}{\sqrt{10}}
\]

\pause

\[
3.228 < \mu < 4.372
\]
\end{frame}

\begin{frame}{Exercise 4}
\protect\hypertarget{exercise-4}{}
The voltage of a power supply plant is believed to follow a normal
distribution. A random sample of 18 observations is taken from this
plant, presented below:

\bigskip
\begin{table}[h!]
\footnotesize
\centering
\begin{tabular}{ccccccccc}
\toprule
10.85 & 11.40 & 10.81 & 10.24 & 10.23 & 9.49 & 9.89 & 10.11 & 10.57 \\
11.21 & 10.10 & 11.22 & 10.31 & 11.24 & 9.51 & 10.52 & 9.92 & 8.33  \\
\bottomrule
\end{tabular}
\end{table}
\medskip

To avoid any need for calculation, consider that the mean and standard
deviation of the above sample is 10.3306 and 0.76672. The primary
concern is the average voltage of a particular power supply plant,
believed to be 10 volts.
\end{frame}

\begin{frame}{Exercise 4 (Cont'd)}
\protect\hypertarget{exercise-4-contd}{}
\begin{enumerate}
\item
  Calculate the 95\% confidence interval for the average voltage for
  this power supply unit. Is the belief that \(\mu = 10\) reasonable?

  \pause

  Since the population variance is unknown and the population
  distribution is normally distributed, we can use Student's t
  distribution to construct the confidence interval. To find confidence
  interval, we first need to find \(t_{\alpha/2}\). We want to construct
  a 95\% confidence interval. Therefore, the confidence level is
  \(1-\alpha = 0.95\), and \(\alpha = 0.05\). Considering that we have a
  two-sided confidence interval, the error term in each tail must be
  \(\alpha/2 = 0.025\). Thus, \(\Pr(T < \alpha) = 0.975\). To find
  suitable \(t_{\alpha/2} = \alpha\), search for the probability of
  0.975 among the numbers (probabilities) with a degree of freedom of
  \(n-1 = 18-1 = 17\) in the Student's t distribution table. We find
  that \(\Pr(T < 2.110) = 0.975\). Therefore, we have
  \(t_{\alpha/2} = 2.110\).
\end{enumerate}
\end{frame}

\begin{frame}{Exercise 4 (Cont'd)}
\protect\hypertarget{exercise-4-contd-1}{}
\pause

Construct the confidence interval and replace \(\bar{x} = 10.3306\),
\(t_{\alpha/2} = 2.11\), \(n = 18\), and \(s = 0.76672\) as follows:

\pause

\[
\bar{x} - t_{\alpha/2} \times \frac{s}{\sqrt{n}} < \mu < \bar{x} + t_{\alpha/2} \times \frac{s}{\sqrt{n}}
\]

\pause

\[
10.3306 - 2.11 \times \frac{0.76672}{\sqrt{18}} < \mu < 10.3306 + 2.11 \times \frac{0.76672}{\sqrt{18}}
\]

\pause

\[
9.9493 < \mu < 10.7119
\]
\end{frame}

\begin{frame}{Exercise 4 (Cont'd)}
\protect\hypertarget{exercise-4-contd-2}{}
\pause

This means that if we were to calculate many 95\% confidence intervals
using the same method, about 95\% of them would contain the true mean.
Since \(\mu = 10\) lies within this interval, the belief that the
population mean is 10 is plausible based on this sample.
\end{frame}

\begin{frame}{Exercise 4 (Cont'd)}
\protect\hypertarget{exercise-4-contd-3}{}
\begin{enumerate}
\setcounter{enumi}{1}
\item
  Calculate the 99\% confidence interval

  \pause

  We want to construct a 99\% confidence interval. Therefore, the
  confidence level is \(1-\alpha = 0.99\), and \(\alpha = 0.01\).
  Considering that we have a two-sided confidence interval, the error
  term in each tail must be \(\alpha/2 = 0.005\). Thus,
  \(\Pr(T < \alpha) = 0.995\). To find suitable
  \(t_{\alpha/2} = \alpha\), search for the probability of 0.995 among
  the numbers (probabilities) with a degree of freedom of
  \(n-1 = 18-1 = 17\) in the Student's t distribution table. We find
  that \(\Pr(T < 2.898) = 0.995\). Therefore, we have
  \(t_{\alpha/2} = 2.898\).
\end{enumerate}
\end{frame}

\begin{frame}{Exercise 4 (Cont'd)}
\protect\hypertarget{exercise-4-contd-4}{}
Now construct the confidence interval as follows:

\pause

\[
\bar{X} - t_{\alpha/2} \times \frac{s}{\sqrt{n}} < \mu < \bar{X} + t_{\alpha/2} \times \frac{s}{\sqrt{n}}
\]

\pause

\[
10.3306 - 2.898 \times \frac{0.76672}{\sqrt{18}} < \mu < 10.3306 + 2.898 \times \frac{0.76672}{\sqrt{18}}
\]

\pause

\[
9.8068 < \mu < 10.8543
\]

Note that by increasing the confidence level, the confidence interval
becomes wider.
\end{frame}

\begin{frame}{Exercise 5}
\protect\hypertarget{exercise-5}{}
The following P/E ratios were obtained for a random sample of 8 stocks
taken from the \textbf{Wall Street Journal}.

\bigskip
\begin{table}[h!]
\centering
\begin{tabular}{cccccccc}
\toprule
5 & 7 & 9 & 10 & 14 & 23 & 20 & 15 \\
\bottomrule
\end{tabular}
\end{table}
\bigskip

Calculate:

\begin{enumerate}
\item
  An estimator for the mean P/E ratio on the New York Stock Exchange.
  State any assumptions that are necessary for your work.

  \pause

  We can estimate the population mean using the sample mean, which in
  this case is 12.875
\end{enumerate}
\end{frame}

\begin{frame}{Exercise 5 (Cont'd)}
\protect\hypertarget{exercise-5-contd}{}
\begin{enumerate}
\setcounter{enumi}{1}
\item
  A 95\% confidence interval for the mean P/E ratio on the New York
  Stock Exchange. State any assumptions that are necessary for your
  work.

  \pause

  Since the population standard deviation is unknown, population
  distribution is unknown, and the sample size is small. We will assume
  that the population distribution is normally distributed so that we
  can use Student's t distribution to construct the confidence interval.
  To find the confidence interval, we first need to find
  \(t_{\alpha/2}\). We want to construct a 95\% confidence interval.
  Therefore, the confidence level is \(1-\alpha = 0.95\), and
  \(\alpha = 0.05\). Considering that we have a two-sided confidence
  interval, the error term in each tail must be \(\alpha/2 = 0.025\).
  Thus, \(\Pr(T < \alpha) = 0.975\).
\end{enumerate}
\end{frame}

\begin{frame}{Exercise 5 (Cont'd)}
\protect\hypertarget{exercise-5-contd-1}{}
\pause

To find suitable \(t_{\alpha/2} = \alpha\), search for the probability
of 0.975 among the numbers (probabilities) with a degree of freedom of
\(n-1 = 8-1 = 7\) in the Student's t distribution table. We find that
\(\Pr(T < 2.364) = 0.975\). Therefore, we have \(t_{\alpha/2} = 2.364\).

\pause

Construct the confidence interval and replace \(\bar{x} = 12.875\),
\(t_{\alpha/2} = 2.364\), \(n = 8\), and \(s = 6.31\) as follows:
\end{frame}

\begin{frame}{Exercise 5 (Cont'd)}
\protect\hypertarget{exercise-5-contd-2}{}
\[
\bar{x} - t_{\alpha/2} \times \frac{s}{\sqrt{n}} < \mu < \bar{x} + t_{\alpha/2} \times \frac{s}{\sqrt{n}}
\]

\pause

\[
12.875 - 2.364 \times \frac{6.31}{\sqrt{8}} < \mu < 12.875 + 2.364 \times \frac{6.31}{\sqrt{8}}
\]

\pause

\[
7.598 < \mu < 18.152
\]
\end{frame}



\end{document}
